% status: ZERO
% Measurement anchor: research/measurements/2026-09-24-m1-part4.md.
\chapter{When Speculation Loses}\label{ch:when-speculation-loses}

\begin{ChapterIntent}
Speculative decoding is not a free acceleration layer. It exchanges extra proposal,
verification, KV and scheduler work for fewer serial target iterations. This chapter
derives the break-even condition, shows why acceptance rate alone is an unsafe control
signal, connects speculation to the roofline, concurrency, context length, paged KV,
continuous batching and SLO goodput, and builds a controller that is allowed to choose
the most important speculative configuration: off.
\end{ChapterIntent}

\OpeningSection{The drafts were useful and the server still became slower}

On the book's Apple M1 measurement, an 8B target decoded at 9.84--10.57 tokens/s on
the quote workload and 10.08--10.29 tokens/s on free writing. Pairing it with a
quantized 3B drafter reduced those ranges to 5.23--5.35 and 4.34--5.57 tokens/s.
The target accepted 98 of 208 drafts on quote and 106 of 231 on free writing:
46--47 percent acceptance. A one-token drafter pass cost 53.4 ms versus 140.7 ms
for the target, or 0.38 target steps
(measured; \texttt{research/measurements/2026-09-24-m1-part4.md}).

A second experiment removes proposal quality as the explanation. Prompt lookup on a
quotation accepted every proposed token. One request improved by about 4.05 times on
the paired server; four simultaneous requests improved by only 1.99--2.03 times.
Acceptance remained 100 percent. What disappeared was idle execution capacity.

The production question is therefore not ``is the drafter accurate?'' It is
\Callout{Production question}{Does speculative execution reduce resource-time
per useful committed token under the current workload and SLO?}

\Foundations{Saved target iterations must pay for all speculative work}{
Count only committed output. Proposals, verified tree nodes and rolled-back KV are
work, not progress. Compare against the best ordinary decode policy at the same
request mix, output semantics and SLO. A speculative policy can improve one request's
latency while reducing fleet goodput by consuming batch slots, KV capacity or compute
that other requests could use. Chapter~\ref{ch:speculative-decoding}
derives the accepted-token count; Chapter~\ref{ch:continuous-batching}
explains why concurrent requests compete for iteration budgets. Put those
two ideas together: the denominator is the time and capacity spent by
the whole policy, not just the target verification kernel. Include the
drafter, rejected positions and control work before claiming a speedup.
}

\NapkinSection{Napkin math: the general break-even inequality}

For an ordinary batch $B$, let target decode step time be $t_1(B,L)$ at representative
context length $L$. Let a speculative iteration propose/verify $N$ extra nodes per
request, commit $A+1$ tokens per request on average, spend draft time $t_d$, target
verification time $t_v$, controller time $t_o$, and amortized interference cost $t_i$.
Speculation wins on raw committed-token time when
\begin{equation}
\frac{t_d+t_v+t_o+t_i}{B(A+1)}
<
\frac{t_1(B,L)}{B}.
\label{eq:general-break-even}
\end{equation}
Equivalently,
\begin{equation}
A+1 >
\frac{t_d+t_v+t_o+t_i}{t_1(B,L)}.
\label{eq:accepted-break-even}
\end{equation}
The right-hand side is the \emph{cost ratio}. It is not constant: $t_v$ changes with
batch, tree width and context; $t_i$ changes with queue pressure and KV occupancy.

For a linear drafter of length $\gamma$, per-draft cost $c$ measured in ordinary target
steps, and verification ratio
\[
v(B,L,\gamma)=\frac{t_v(B,L,\gamma)}{t_1(B,L)},
\]
a useful approximation is
\begin{equation}
S \approx \frac{\tau}{\gamma c+v+o+i},
\qquad
\text{win iff }\tau>\gamma c+v+o+i ,
\label{eq:linear-break-even}
\end{equation}
where $\tau=A+1$ and $o,i$ are normalized overhead/interference.

\section{Acceptance is numerator quality, not profitability}

A controller that enables speculation when acceptance exceeds, say, 70 percent ignores
the denominator. Consider $\gamma=4$ and an average $\tau=4$ committed tokens. If
$c=.05$, $v=1.2$ and overhead is $.1$, normalized cost is $1.5$ and the speedup estimate
is $2.67$. If $c=.4$, $v=2.5$ and overhead is $.2$, cost is $4.3$ and the same accepted
work loses.

Conversely, modest acceptance can pay when proposals are nearly free and verification
runs in otherwise idle arithmetic. The controller must estimate both accepted progress
and marginal resource cost.

\section{From roofline to break-even surface}

Autoregressive decode at small batch is commonly weight-memory-bandwidth dominated:
one token cannot use all matrix arithmetic while model weights must still be streamed.
Speculation increases tokens evaluated per target pass and can fill idle arithmetic.

Let $t(n,L)$ be target step time for $n$ token positions at context $L$. Ignoring draft
cost for a moment, linear speculation with $\gamma$ drafts wins when
\begin{equation}
\frac{t(B(\gamma+1),L)}{t(B,L)} < \tau.
\label{eq:batch-break-even}
\end{equation}
If both points lie on a flat bandwidth-dominated region, the ratio approaches one.
When verification crosses into compute saturation, it rises. Thus break-even is a
surface over
\[
(B,L,\gamma,\tau,\text{model},\text{dtype},\text{kernel},\text{hardware}),
\]
not one magic batch size.

\EngineFigure{ch24-break-even-surface}{Decision view of the break-even condition, not a measured surface plot. Batch, context, draft cost and accepted path length jointly determine whether speculation adds useful goodput.}{fig:spec-break-even}

\section{The M1 verification curve shows the ridge directly}

The measured 3B target needed 68.3 ms for one token and 157.5 ms for 16 tokens:
16 positions cost only 2.30 single-token steps. The 8B target needed 140.7 ms for one
and 379.6 ms for 16, a ratio of 2.70. At 32 positions, ratios were only 2.33 and 2.78.

Those are precisely the circumstances in which verifying several positions can be
profitable: work grows much faster than elapsed time. But the curve is not flat.
Speculation that assumes ``16 verified tokens cost one decode'' would overestimate
speedup substantially on this machine.

The non-monotonic small-$n$ measurements (for example 3 versus 4 positions) also warn
against fitting a simplistic line to kernel dispatch. Measure the actual shapes the
serving engine uses.

\section{Concurrency spends the same idle arithmetic}

Without speculation, continuous batching increases $B$ by adding independent requests.
With speculation, intra-request parallelism increases target work by adding candidate
positions. Both attempt to occupy the same hardware.

At low arrival rate, spending spare compute on one request can sharply reduce ITL. At
high arrival rate, those speculative positions compete with queued requests. The
opportunity cost $t_i$ in equation~\eqref{eq:general-break-even} becomes real even if
the speculative request itself finishes faster.

The M1 quote experiment isolates this effect: 100-percent acceptance, same lookup
mechanism, roughly 4.05 times faster for one request but about 2.0 times for four.
The correct interpretation is not that speculation ``stops working'' at concurrency
four. Its marginal value fell because ordinary batching became a better consumer of
the machine.

\section{Large batch is not universally hostile: context length changes the bottleneck}

A simplistic rule ``disable speculation at large batch'' is also wrong. Long-context
decode reads a large KV history. As context grows, attention/KV traffic can return the
system to a memory-bound regime even at large batch. MagicDec analyzes this regime and
uses a sparse-KV drafter to reduce draft cost; its authors report gains for large-batch,
long-context serving \cite{sadhukhan2024magicdec}.

The systems lesson is broader than that method: the ridge moves with $L$. A controller
trained only on short-context profiling can make the wrong decision for long-context
requests.

\section{Draft cost has three parts}

Do not model draft cost as parameter count alone:
\[
t_d=t_{\rm weights}+t_{\rm draftKV}+t_{\rm coordination}.
\]
A separate drafter reads weights, maintains its own KV and may require tokenizer
translation or device communication. An EAGLE-style drafter may reduce weight cost but
require target-feature movement. A parallel drafter may reduce serial steps but perform
more one-shot computation.

The M1 3B drafter had tensor bytes 0.37 of the 8B target and measured one-token time
0.38 of the target. With several serial drafts, that term alone can exceed a target
decode step before verification begins.

\section{Memory capacity has an opportunity cost}

Speculation consumes memory through drafter weights, drafter KV, tentative target KV,
tree metadata and potentially larger workspace. Let ordinary serving fit $B_{\max}$
requests but speculation reduce it to $B'_{\max}$. Then even a per-request latency
win can lower fleet throughput if
\[
G_{\rm spec}(B'_{\max}) < G_{\rm base}(B_{\max}).
\]

For the measured 8B/3B pair, drafter tensor bytes were about 1.92 GB against 5.22 GB
target tensor bytes---roughly 37 percent additional target-weight size before draft KV
and runtime buffers. On capacity-constrained accelerators this can be more damaging
than draft compute.

\section{Tree speculation magnifies wasted verification}

For a chain, rejected work after the first mismatch is easy to visualize. For a tree,
many verified nodes are intentionally incompatible alternatives. Define
\[
\eta_{\rm tree}=\frac{A+1}{N_v},
\]
where $N_v$ is verified tree nodes. Width can increase accepted path length while
lowering this efficiency.

Under low load, a wide tree may be worthwhile because the target pass has unused
parallel capacity. Under high load, those losing branches compete directly with other
requests. Dynamic tree size should therefore be load-aware, not only confidence-aware.

\section{KV pressure changes the scheduler before OOM}

Suppose a speculative tree reserves $R_s$ tentative blocks and ordinary decode would
reserve $R_b$. The relevant cost is not only whether $R_s$ fits. Extra reservations can
force another request to wait, prevent admission, trigger preemption, or reduce prefix
cache residency.

Define a KV pressure ratio
\[
\rho_{\rm KV}=\frac{\text{allocated + reserved blocks}}{\text{usable blocks}}.
\]
As $\rho_{\rm KV}$ approaches a policy threshold, shrink tree/draft budgets before
hard allocation failure. Speculation should be one of the first optional consumers to
yield memory.

\section{Chunked prefill and speculation share an iteration budget}

A server may have queued prompt chunks, ordinary decodes and speculative verification
at the same time. If the iteration budget is $T$, every speculative node can displace
prefill progress.

This creates an SLO trade: speculation may improve ITL for already-running requests
while increasing TTFT for newcomers. A goodput controller must include both objectives
or enforce separate budgets/reservations.

One safe hierarchy is:
\begin{enumerate}
\item Meet mandatory decode and SLO work.
\item Preserve minimum fair prefill progress.
\item Spend remaining capacity on profitable speculative work.
\end{enumerate}
That ordering is policy, not a universal optimum, but it makes speculation explicitly
preemptible by required service work.

\section{Sampling changes proposal agreement}

Higher temperature does not mechanically imply a fixed acceptance drop; it changes the
target distribution and therefore proposal agreement according to the local entropy
and exact speculative algorithm.

The M1 prompt-lookup quote remained perfectly accepted at temperature 0.8 because
copying the quotation was highly predictable. Edit acceptance changed only slightly.
Free/summarize tasks produced no prompt-lookup drafts and showed ordinary run-to-run
noise. The correct controller feature is observed accepted progress for the current
workload, not temperature alone.

\section{Structured output can help or hurt}

A grammar can prune illegal draft branches early, raising useful proposal density.
But grammar state must also be advanced for every candidate path that is actually
verified. Complex branch-local state adds controller cost and memory.

If constraints make the next token nearly deterministic, speculation may become
excellent. If constraint checks dominate a tiny target, the opposite can happen.
Measure the integrated sampler/grammar path rather than benchmarking raw target logits.

\section{Streaming exposes burst latency}

Speculation emits several committed tokens after one verification boundary. Average
time per token can improve while users observe a longer gap followed by a burst.

For interactive service track both:
\[
\mathrm{ITL}_{\rm commit}
\quad\text{and}\quad
\mathrm{gap}_{\rm verify}.
\]
A configuration that maximizes tokens/s can violate a maximum-gap SLO. The controller
must be allowed to shorten $\gamma$ even when throughput would prefer a longer draft.

\section{Goodput is the serving objective}

Throughput counts all completed tokens. SLO goodput counts useful completions that meet
the service contract. One form is
\[
G_{\rm SLO}=
\frac{\sum_r y_r\,
\mathbf{1}[\mathrm{TTFT}_r\le T_r^{\max}]
\mathbf{1}[\mathrm{ITL}_{r,p}\le I_r^{\max}\ \forall p]}
{\Delta t},
\]
where $y_r$ is committed output.

The exact business metric may use percentiles rather than all-token constraints. The
principle is stable: speculative work that increases raw throughput but causes more
requests to miss latency objectives is not an optimization.

\section{Counterfactual measurement is the hard part}

When speculation is enabled, the engine observes speculative cost but not the baseline
time the same batch would have taken at that instant. When disabled, it cannot directly
observe current acceptance.

A controller therefore needs a model, periodic exploration, paired profiling, or
shadow estimates. TurboSpec is an example of treating speculation as a closed-loop
goodput-control problem rather than a static flag \cite{liu2024turbospec}.

Avoid an estimator trained only on successful speculative traffic: disabling
speculation changes batching, which changes the very verification-cost curve being
predicted.

\section{A practical controller}

Let configurations be
\[
\mathcal C=\{0,1,2,4,8,\ldots\},
\]
where zero disables speculation and positive values represent draft/tree budgets.
For each configuration maintain exponentially weighted estimates of committed tokens,
draft time, verification time, iteration time, KV reservations and SLO misses.

For configuration $c$, predict
\[
\widehat G(c)=
\frac{\widehat{\text{committed useful tokens}}}
{\widehat{\text{resource time}}}
-\lambda_{\rm SLO}\widehat P_{\rm miss}
-\lambda_{\rm KV}\widehat R_{\rm KV}.
\]
Choose the best feasible $c$, but add hysteresis:
\[
\widehat G(c_{\rm new}) >
(1+\epsilon)\widehat G(c_{\rm current})
\]
for $m$ windows before switching. Hysteresis prevents measurement noise from causing
oscillation.

\EngineFigure{ch24-controller}{The speculation controller observes committed progress, execution cost, queue/KV pressure and SLO misses; it selects a draft/tree budget, including zero, and uses hysteresis before switching.}{fig:spec-controller}

\section{Control timescale matters}

Changing draft length every token reacts to noise. Waiting minutes reacts too slowly to
traffic shifts. A useful hierarchy has:
\begin{itemize}
\item per-iteration hard safety limits from token/KV budgets;
\item short-window adaptation of draft/tree size;
\item slower recalibration of cost models as hardware/load regimes change.
\end{itemize}
Request-level hints can initialize policy---for example long-context versus short---but
observed performance should dominate.

\section{Avoid positive feedback under overload}

A dangerous controller sees falling acceptance, increases tree width to find more
alternatives, lengthens iterations, grows the queue, and then widens again because
per-request progress is poor.

Overload logic must move the other direction. As queue delay, KV pressure or SLO misses
rise, cap optional speculative work. Admission/scheduling pressure is a reason to
become conservative even if isolated-request speculation would win.

\section{Per-tenant and per-model control}

Acceptance and cost differ by model, adapter, prompt domain, sampling settings and
tenant workload. One global EWMA can average incompatible regimes.

Key statistics by a stable configuration signature such as
\[
H(\text{target},\text{drafter},\text{adapter},\text{sampling class},
\text{context bucket},\text{hardware/backend}).
\]
Do not include sensitive prompt text. Bucketing must have fallback priors so sparse
classes do not overfit.

\section{Benchmark the policy, not only the mechanism}

An offline speculative benchmark with fixed batch and fixed $\gamma$ answers a kernel
question. A serving benchmark must vary arrivals, contexts, output lengths and
acceptance regimes.

Report baseline and speculative TTFT, ITL, throughput, SLO goodput, acceptance/accepted
length, verified nodes, draft/verify/controller time, KV high-water mark, preemptions,
prefix-cache effects and configuration-switch count.

Also report \emph{off-policy regret}: how often the controller chose speculation when
the measured/estimated baseline would have been better, and vice versa.

\section{Failure injection}

A production controller should be tested against:
\begin{enumerate}
\item sudden shift from copying to creative generation;
\item concurrency ramp from one to saturation;
\item long-context burst that changes the memory/compute bottleneck;
\item KV occupancy spike;
\item drafter slowdown or device transfer regression;
\item noisy timing samples;
\item SLO tightening for a priority tenant.
\end{enumerate}
The safe failure mode is generally less speculation, not more.

\EngineFigure{sys-speculation-choice}{A break-even flowchart for speculation. Compare time per committed token under the same quality, load and latency contract. Draft, verify, rollback and queue costs all belong in the measurement.}{fig:sys-speculation-choice}

\LedgerSection
\begin{ConstraintLedger}
Memory bandwidth & trades & Speculation exploits otherwise idle arithmetic during bandwidth-bound target steps; long KV histories can restore this opportunity at large batch.\\
Memory capacity & spends & Drafter weights/KV, tentative target KV, tree state and workspace can reduce batch capacity or prefix-cache residency.\\
Compute & spends & Rejected chain suffixes and losing tree branches become direct opportunity cost once ordinary batching saturates arithmetic.\\
Latency & trades & Fewer serial target iterations can reduce ITL, but verification bursts and queue interference can worsen tail gaps and TTFT.\\
Correctness & requires & Compare equal decoding semantics; never count proposals as output and never publish losing speculative state.\\
Cost & decides & Enable speculation only when committed SLO-goodput per resource improves over the counterfactual baseline.\\
\end{ConstraintLedger}

\LabSection{a controller that is allowed to turn speculation off}

The Rust lab in \texttt{code/labs/ch24-break-even/} is a deterministic trace simulator.
Each observation supplies baseline step time, speculative draft/verify/overhead time,
committed tokens, KV pressure and SLO miss state. It computes derived speedup,
updates an EWMA score for each draft budget, and applies hysteresis before switching.
The trace values are illustrative inputs, not timings measured by this lab.
For positive scores the hysteresis margin is proportional to the current
score; for negative penalized scores it uses the current score's magnitude
so a small improvement still cannot trigger an immediate switch.

The independent oracle evaluates the break-even inequality directly from raw trace
fields. Tests cover a profitable low-load regime, an expensive-drafter loss, a
100-percent-acceptance case whose value falls with concurrency, KV-pressure disabling,
and hysteresis against oscillation. This lab does not claim to reproduce GPU timings.

\HappenedSection{the same technique crossed both sides of break-even}

The M1 record contains both outcomes. Prompt lookup on repeated text was strongly
profitable, especially at one request. A mismatched 3B drafter on the 8B reasoning
target was strongly unprofitable despite accepting nearly half its proposals. Four-way
concurrency reduced the benefit of perfect prompt-lookup proposals without changing
their acceptance.

Research systems have consequently moved toward workload-aware control. TurboSpec
frames the problem as closed-loop goodput optimization \cite{liu2024turbospec}.
MagicDec shows why long context can reopen a speculative opportunity at large batch:
KV traffic changes the bottleneck and a sparse-KV drafter attacks draft cost
\cite{sadhukhan2024magicdec}.

\FrontierSection{2026-10-05}

The frontier is no longer a binary argument over whether speculative decoding ``works.''
It is control over a multi-dimensional operating region. Proposal architecture changes
acceptance and draft cost; verification kernels change $v$; continuous batching changes
opportunity cost; context changes the roofline; KV managers change capacity cost; SLOs
change the objective.

A mature inference engine should therefore expose speculation as a policy with
observable economics:
\[
\boxed{
\text{speculate only while marginal committed SLO-goodput exceeds marginal cost.}
}
\]
The next chapters return to KV architecture, where context representation changes both
ordinary decode cost and the break-even surface derived here.

\ExercisesSection
\begin{enumerate}
\item \emph{Derivation.} With $t(n)=\max(t_0,an)$, zero draft cost and linear draft
length $\gamma$, derive the regions of $B$ for which equation~\eqref{eq:batch-break-even}
is bandwidth-bound/bandwidth-bound, bandwidth-bound/compute-bound, or
compute-bound/compute-bound.
\item \emph{Measured arithmetic.} Using the M1 3B values 68.3 ms at one token and
157.5 ms at 16, compute verification ratio $v$. What minimum $\tau$ is required if
draft+controller cost is 0.5 target steps?
\item \emph{Capacity.} A drafter reduces maximum resident batch from 64 to 48 but makes
each resident request 1.25 times faster. Ignoring other effects, does aggregate capacity
improve?
\item \emph{Controller.} Add separate context-length buckets to the lab and construct a
trace where one global estimator chooses the wrong policy.
\item \emph{SLO.} Construct two policies with equal raw throughput but different TTFT
and ITL distributions; show how SLO goodput reverses the ranking.
\item \emph{Falsification.} Produce a deterministic trace with 100-percent acceptance
where speculation loses. Identify the denominator term responsible.
\end{enumerate}

\Needspace{18\baselineskip}
\section{Worked problems}

\Needspace{12\baselineskip}
\subsection{Piecewise ridge}
\textbf{Problem.} \emph{Derivation.} With $t(n)=\max(t_0,an)$, zero draft cost and linear draft
length $\gamma$, derive the regions of $B$ for which equation~\eqref{eq:batch-break-even}
is bandwidth-bound/bandwidth-bound, bandwidth-bound/compute-bound, or
compute-bound/compute-bound.

\textbf{Solution.} With $t(n)=\max(t_0,an)$, let $n^*=t_0/a$.
If $B(\gamma+1)\le n^*$, verification and baseline are both flat and the ratio is 1.
If $B<n^*<B(\gamma+1)$, the ratio is $aB(\gamma+1)/t_0$.
If $B\ge n^*$, the ratio is $\gamma+1$. Compare each with $\tau$.

\Needspace{12\baselineskip}
\subsection{M1}
\textbf{Problem.} \emph{Measured arithmetic.} Using the M1 3B values 68.3 ms at one token and
157.5 ms at 16, compute verification ratio $v$. What minimum $\tau$ is required if
draft+controller cost is 0.5 target steps?

\textbf{Solution.} $v=157.5/68.3\approx2.31$. With another 0.5 normalized step of
draft/controller work, break-even requires $\tau>2.81$ committed tokens per iteration.

\Needspace{12\baselineskip}
\subsection{Capacity}
\textbf{Problem.} \emph{Capacity.} A drafter reduces maximum resident batch from 64 to 48 but makes
each resident request 1.25 times faster. Ignoring other effects, does aggregate capacity
improve?

\textbf{Solution.} Aggregate relative capacity is $(48/64)\times1.25=0.9375$.
Under the stated simplification, it falls by 6.25 percent.

\Needspace{12\baselineskip}
\subsection{Perfect acceptance can lose}
\textbf{Problem.} \emph{Falsification.} Produce a deterministic trace with 100-percent acceptance
where speculation loses. Identify the denominator term responsible.

\textbf{Solution.} Let $\gamma=4$, $\tau=5$, draft cost be 2.0
target steps and verification cost 3.5 steps. Normalized total is 5.5, so
$S=5/5.5<1$ despite perfect acceptance.

